%%=============================================================
%%  Stage 11 Expanded Section Notes — LaTeX scaffolding
%%  Source: PAPER_SECTION_NOTES_EXPANDED.md (Stage 11, prose form)
%%  Template target: memoir-class thesis
%%    abstract → 01_chapter1 → 02_chapter2 → 03_chapter3
%%               → 04_chapter4 → 05_chapter5 → appendix
%%
%%  THIS IS DRAFTING SCAFFOLDING, NOT FINAL THESIS PROSE.
%%  Every numeric value is tied to the immutable locators in
%%  PAPER_SECTION_NOTES.md and the committed origin/results CSVs.
%%  Evidence labels: [DIRECT] [DERIVED] [INTERPRETATION]
%%                   [UNRESOLVED] [SUPERSEDED]
%%  Research-status tags follow 00_PROTOCOL.md §4.
%%  No phase may be claimed PRE_REGISTERED_CONFIRMATORY
%%  (formal post-pivot preregistration is NOT FOUND, GAP-021).
%%=============================================================

\documentclass[12pt, a4paper, oneside]{memoir}

%% ----- packages -----
\usepackage{lmodern}
\usepackage[T1]{fontenc}
\usepackage[utf8]{inputenc}
\usepackage{microtype}
\usepackage{amsmath}
\usepackage{booktabs}
\usepackage{hyperref}
\usepackage{xcolor}
\usepackage{enumitem}
\usepackage{csquotes}

\hypersetup{colorlinks=true, linkcolor=blue, citecolor=blue, urlcolor=blue}

%% ----- evidence-label commands -----
\newcommand{\evDirect}{\textcolor{blue}{[DIRECT]}}
\newcommand{\evDerived}{\textcolor{teal}{[DERIVED]}}
\newcommand{\evInterp}{\textcolor{violet}{[INTERPRETATION]}}
\newcommand{\evUnresolved}{\textcolor{red}{[UNRESOLVED]}}
\newcommand{\evSuperseded}{\textcolor{gray}{[SUPERSEDED]}}

%% research-status
\newcommand{\rsPHE}{\texttt{POST\_HOC\_EXPLORATORY}}
\newcommand{\rsEV}{\texttt{ENGINEERING\_VALIDATION}}

%% convenience
\newcommand{\qlT}{\texttt{ql\_true}}
\newcommand{\qlF}{\texttt{ql\_false}}
\newcommand{\wsRD}{\texttt{ws:responseDelay}}

%% ----- document -----
\title{Stage 11 Expanded Section Notes\\[0.5ex]
       \large LaTeX Scaffolding for Thesis Manuscript}
\author{MT-Esra}
\date{2026-06-18}

\begin{document}
\maketitle
\tableofcontents
\clearpage

%%=============================================================
\chapter*{Reading and Labelling Conventions}
\addcontentsline{toc}{chapter}{Reading and Labelling Conventions}
%%=============================================================

Evidence labels \evDirect, \evDerived, \evInterp, \evUnresolved, and
\evSuperseded\ carry the same meaning as in \texttt{00\_PROTOCOL.md} §3.
Research-status tags (\rsPHE, \rsEV, etc.) follow \texttt{00\_PROTOCOL.md}
§4. No phase may be written up as
\texttt{PRE\_REGISTERED\_CONFIRMATORY}: formal post-pivot preregistration is
\texttt{NOT FOUND} (GAP-021 / S9T-019 / P3T-035).

Numbers stated in these notes are either \emph{deterministic} (recomputable
from raw run data) or \emph{direct analysis outputs} (bootstrap confidence
intervals and bootstrap $p$-values that were \textbf{not} independently
resampled in this environment).  The prose flags which is which wherever a
confidence interval or bootstrap $p$-value appears, because the Stage~9 audit
kept those two categories explicitly separate.

%%=============================================================
\chapter{Abstract}
\label{ch:abstract}
% Target file: content/abstract.tex
% Write the abstract last, after the chapters are settled.
%%=============================================================

The abstract follows a four-move structure: problem, method, principal
findings, strongest limitation.

\section{Problem move \evDirect}

The thesis studies whether semantic / knowledge-graph (KG) and physics
knowledge measurably helps a reinforcement-learning agent control small
simulated building labs, across three increasingly demanding situations:
(i)~learning a clean lab whose physics already matches the KG;
(ii)~reacting when a component the agent was trained on becomes defective;
and
(iii)~learning the temporal response behaviour of actuators and feeding that
knowledge back into the KG\@.
These three situations come directly from the advisor's requirement note and
motivate the Phase~1 / Phase~2 / Phase~3 framing.

\section{Method move \evDirect}

Three mechanisms are evaluated on a JaCaMo (Jason BDI + CArtAgO) multi-agent
stack with a tabular Q-learner, a KG/stereotype reasoner, and Node-RED
illuminance-lab simulators:
\begin{enumerate}
  \item clean-lab KG/physics priming of the Q-learner;
  \item an expected-versus-actual fault detector that blacklists a defective
        component and warm-restarts learning; and
  \item a probing routine that estimates per-actuator response delays, writes
        them back to the KG as \wsRD, and lets a deadline-aware BDI planner
        use them.
\end{enumerate}

\section{Principal findings}

\begin{description}
  \item[Phase~1 \evDerived, \rsPHE]
    In the clean lab2 profile, the KG/physics-primed arm (\qlT) beats the
    no-knowledge arm (\qlF) on goal-reaching AUC, average steps, redundant
    actions, and goal rate; the improvement also replicates on fresh seeds.
    Lab3 is more nuanced: its reward signal improves clearly while its
    \emph{binary} goal metric does not.

  \item[Phase~2 \evInterp, \rsEV; \rsPHE]
    The detector achieves complete detection recall and correct component
    attribution across every injected-fault cell tested; the
    \emph{recovery-speed} advantage of the KG arm is narrow --- faster for
    \texttt{lab3\_f1dead} but not for \texttt{lab3\_f1inv}.

  \item[Phase~3 \evInterp, \rsEV; \rsPHE]
    The agent learns the blind's delay to within ${\approx}2\,\%$ of ground
    truth and the informed planner meets every tight deadline where the
    delay-blind planner meets none, at the cost of higher energy use.
\end{description}

\section{Limitation move \evUnresolved}

State plainly that:
this is controlled simulator/lab evidence, not real-building deployment;
no phase was formally preregistered;
some statistical fields (Phase~1/2 bootstrap CIs and bootstrap $p$-values)
are reported as original analysis outputs and were not independently
re-resampled; and
exact workflow-dispatch payloads and artifact-ZIP byte identity remain
unverified.

%%=============================================================
\chapter{Introduction}
\label{ch:introduction}
% Target file: content/01_chapter1.tex
%%=============================================================

\section{Problem context \evDirect}
\label{sec:intro-problem}

Open from the advisor's anchoring question: in a \emph{clean} environment,
does a Q-learning agent that is given prior physics/KG knowledge do better
than one that starts tabula rasa?  The expected gains were specifically
\emph{time-to-goal} and \emph{avoidance of redundant actions}, with final
task success expected to be roughly comparable between the two.  This
expectation, not a result, motivates the work and frames the first research
question.  Establish here that the agent operates on small
illuminance-control \enquote{labs}---rooms with lamps, a spotlight, and
motorised blinds---and that a knowledge graph encodes the nominal physics of
those rooms.

\section{Research trajectory \evDirect}
\label{sec:intro-trajectory}

The project deliberately unfolds in three phases of rising difficulty, each
building on the previous learner without rewriting it.

\begin{description}
  \item[Phase~1 --- clean labs.]
    Labs whose true physics is aligned with the KG\@; the question is purely
    whether KG priors accelerate clean learning.

  \item[Phase~2 --- weakness labs.]
    The agent is first trained clean, then dropped into a lab where a
    component it relied on is now broken (dead or inverted).  The question is
    whether the agent can notice, discard the broken component, and relearn
    --- and whether KG knowledge helps it realign faster.

  \item[Phase~3 --- process dynamics.]
    The agent learns the \emph{temporal} behaviour (response delay) of each
    actuator, writes it back into the KG, and uses it to satisfy
    deadline-constrained goals.
\end{description}

The strict-separation design is important: each phase is additive, so
Phase~1 and Phase~2 numerical results stay reproducible even after later
phases are built.

\section{Research questions}
\label{sec:intro-rqs}

None of the three research questions below is preregistered confirmatory
(formal preregistration is \texttt{NOT FOUND}); all three are framed as
exploratory / engineering-validation questions answered against a local
pre-run expectation.

\begin{description}
  \item[RQ1 \evInterp, \rsPHE]
    Does KG/physics priming improve clean-lab Q-learning \emph{speed} and
    \emph{action efficiency} while preserving task success?

  \item[RQ2 \evInterp, \rsEV; \rsPHE]
    Can a clean-trained agent recognise unexpected component behaviour,
    discard the defective component, and relearn --- and does the KG/physics
    prior speed realignment?

  \item[RQ3 \evInterp, \rsEV; \rsPHE]
    Can the agent estimate actuator response delays, materialise them as KG
    facts, and use them to meet time-bounded goals better than a delay-blind
    planner?
\end{description}

For each RQ, the candidate-RQ table in \texttt{RESEARCH\_MODEL.md} records
the expected mechanism and the falsification condition; carry those into the
chapter so each RQ is stated together with what would have counted as a
negative result.

\section{Contributions \evDirect}
\label{sec:intro-contributions}

Four contributions are claimed.

\begin{enumerate}
  \item A \textbf{semantic/KG prior for tabular Q-learning} in clean
        building labs, including a cross-zone structural prior that is kept
        conceptually distinct from generic KG priming.

  \item A \textbf{fault-recognition loop} that compares KG-predicted effects
        against observed effects, blacklists the offending component, and
        warm-restarts learning over the surviving action set.

  \item A \textbf{response-delay learning mechanism} that probes actuators,
        writes \wsRD\ back to the KG, and drives a deadline-aware BDI
        planner.

  \item A \textbf{reproducible evidence layer}---labelled ledgers, immutable
        locators, and a statistical-integrity audit---that makes the
        provenance and the limits of every reported number explicit.
\end{enumerate}

\section{Scope statement \evUnresolved}
\label{sec:intro-scope}

The evidence is from controlled Node-RED simulators and small labs, not
deployed buildings.  Do not claim formal preregistration, do not claim
byte-verified artifact-ZIP identity, and do not present Phase~1/2 bootstrap
intervals as independently reproduced.  These scope caveats are not
incidental---they are the same items that reappear as the Chapter~5
limitations.

%%=============================================================
\chapter{Background and Related Work}
\label{ch:background}
% Target file: content/02_chapter2.tex
%%=============================================================

Keep \emph{foundations} (what the reader must know) separate from
\emph{related work} (how the project compares to alternative research).  All
external sources cited below are in \texttt{BIBLIOGRAPHIC\_LEDGER.csv}
(\texttt{BIB-S10B-001}\ldots\texttt{016}); DOI-only rows are
metadata-verified but not full-text reviewed (see the full-text gap,
GAP-046).

\section{Foundations}
\label{sec:bg-foundations}

\subsection{MAS and BDI \evDirect}

Introduce the BDI agent model and the AgentSpeak(L) programming language
(Rao 1996, \texttt{BIB-S10B-001}) and the JaCaMo multi-agent programming
platform (Boissier et al.\ 2013, \texttt{BIB-S10B-002}).  These are the
implementation foundations: the project's agents are Jason/AgentSpeak plans
and its environment tools are CArtAgO artifacts (e.g.\ the dynamics agent
and \texttt{DynamicsLearner}).  Treat these as foundations only---JaCaMo
literature does not supply evidence for the project's KG-prior or
delay-learning claims.

\subsection{Q-learning \evDirect}

Present tabular Q-learning via Watkins and Dayan (1992,
\texttt{BIB-S10B-003}) as the RL algorithm the project builds on.
Explicitly note that the Watkins--Dayan convergence theorem must \emph{not}
be used as evidence that the project's runs converged; convergence in these
labs is an empirical matter checked per run.

\subsection{Semantic Web and KG metadata \evDirect}

Introduce SOSA/SSN (\texttt{BIB-S10B-008}, \texttt{BIB-S10B-009}) and Brick
(\texttt{BIB-S10B-010}) as the standard vocabularies for describing sensors,
actuators, and building metadata.  Make the boundary clear: the project uses
a \emph{custom} \texttt{ws:}~stereotype vocabulary (e.g.\
\texttt{building\_2\_slow.ttl}, \texttt{building\_3\_slow.ttl}), so do
\textbf{not} claim SSN/SOSA or Brick compliance or extension.  Note the
source-hygiene correction (a wrong Brick DOI candidate was rejected and the
correct DOI recorded: \texttt{BIB-S10B-REJ-001} $\to$ \texttt{BIB-S10B-010})
if useful.

\section{Related Work}
\label{sec:bg-relwork}

Organise related work into five comparison buckets, each ending with the
explicit difference from this project.  The taxonomy maps onto Table~RW-1.

\subsection{Building-control RL \evDirect}

Compare with deep RL for HVAC control (Wei et al.\ 2017,
\texttt{BIB-S10B-004}), model-free control of thermostatically controlled
loads (Claessens et al.\ 2018, \texttt{BIB-S10B-005}), sparse-observation
direct load control (Ruelens et al.\ 2019, \texttt{BIB-S10B-006}), and MARL
for HVAC (Bayer \& Pruckner 2022, \texttt{BIB-S10B-007}).  These motivate RL
for building control but differ from the project in domain (HVAC/thermal
vs.\ illuminance), in the use of semantic priors, and in the BDI/KG
layering.  State the domain and actuator mismatch explicitly; do not
conflate the project's MAS architecture with multi-agent RL training.

\subsection{KG priors and reward shaping \evDirect/\evInterp}

Use KG-enhanced RL in recommender systems (Zhou et al.\ 2020,
\texttt{BIB-S10B-011}) and the equivalence of potential-based shaping and
Q-value initialisation (Wiewiora 2003, \texttt{BIB-S10B-012}) as
\emph{analogical} and \emph{theoretical} support that structured prior
knowledge can improve sample efficiency and action selection.  Stress that
KGQR is not building automation and Wiewiora is theory, not semantic KG\@;
the project's causal evidence for KG priors must come from the Phase~1
result rows, not from these sources.

\subsection{Physics-informed and model-assisted learning \evDirect}

Compare with PhysQ (Gokhale et al.\ 2022, \texttt{BIB-S10B-013}) and
model-assisted sensor-fault-tolerant HVAC control (Xu et al.\ 2021,
\texttt{BIB-S10B-014}).  These share the motivation that prior physical
knowledge improves building-control learning; the difference is that the
project uses \emph{explicit symbolic RDF/KG priors} and tabular/BDI
mechanisms plus explicit response-delay writeback, not physics-informed
neural representations.

\subsection{Fault and anomaly detection \evDirect}

Compare Phase~2 with transfer learning for HVAC fault detection (Dowling \&
Zhang 2020, \texttt{BIB-S10B-015}) and KG-based industrial root-cause
diagnosis (Chen et al., \texttt{BIB-S10B-016}), plus the model-assisted
fault-tolerant work (\texttt{BIB-S10B-014}).  The project differs by using an
\emph{expected-versus-actual} KG check followed by blacklisting and
relearning, rather than training a transfer classifier or doing industrial
root-cause analysis.

\subsection{Learned temporal and process dynamics \evInterp}

Position Phase~3 against the model-free / sparse-observation control works
(\texttt{BIB-S10B-005}, \texttt{BIB-S10B-006}) and PhysQ
(\texttt{BIB-S10B-013}): building-control RL usually optimises a policy
directly, whereas this project evaluates whether explicit per-actuator
response-delay knowledge can be \emph{learned, written back to the KG, and
consumed by a deadline-aware BDI planner}.  Phrase the novelty as bounded to
the tested slow-lab actuator-delay setting; do not claim the project learns
full building thermal dynamics.

\section{Related-work artefacts \evUnresolved}
\label{sec:bg-artefacts}

Plan Table~RW-1 (rows = the five buckets above; columns = representative
source IDs, method family, mapped project phase, difference from project)
and Figure~RW-1 (a layered map: MAS/BDI execution at the base, KG/stereotype
priors above it, the Q-learning/building-control RL layer, a fault/anomaly
branch, and a learned-response-delay branch, with the project placed at the
intersection of KG priors, the BDI planner, and learned temporal KG
writeback).

Carry the full-text gap note \evUnresolved: DOI-only rows need full-text
review before any technical claim beyond title, venue, and year.

%%=============================================================
\chapter{Contributions and Approach}
\label{ch:contributions}
% Target file: content/03_chapter3.tex
%%=============================================================

Keep this chapter conceptual.  Push file-by-file code mapping to the
appendix; only enough implementation detail to define each mechanism belongs
here.

\section{Approach overview \evInterp}
\label{sec:approach-overview}

The work is one coherent research artefact with three layers of semantic
guidance: semantic priors that \emph{accelerate} clean learning (Phase~1),
semantic expected-versus-actual checks that confer \emph{robustness} under
faults (Phase~2), and learned process dynamics that are \emph{written back}
to the KG and reused for time-bounded control (Phase~3).  The unifying idea
is that a knowledge graph is not only a static prior but also a destination
for newly learned facts.

\section{Phase~1 --- clean-lab ladder and semantic priors \evDirect}
\label{sec:approach-phase1}

Describe the lab ladder (\texttt{lab1}, \texttt{lab2}, \texttt{lab3}) of
increasing state-space size, the two learning arms (\qlF\ no-prior
vs.\ \qlT\ KG-primed), and the benchmark modes (\texttt{rule\_based},
\qlF, \qlT).  The KG/stereotype reasoner supplies action guidance that
belongs to the initialisation / reward-shaping family (consistent with the
Wiewiora framing in Chapter~\ref{ch:background}).

The \textbf{cross-zone structural prior}---which encodes that actuators in
one zone spill light into adjacent zones---is kept conceptually separate from
generic KG priming, because the Phase~1 ablation specifically tests the
\emph{targeted} cross-zone structure against an equal-budget \emph{untargeted}
optimism baseline (\texttt{phase1\_kg\_xzone\_rand},
\texttt{cross\_zone\_bonus=3.0}).

\section{Phase~2 --- fault detection, blacklist, and relearn \evDirect}
\label{sec:approach-phase2}

Present the detect $\to$ attribute $\to$ blacklist $\to$ warm-restart $\to$
re-certify loop.  The agent warm-starts from the clean Phase-1 Q-table, then
on every step compares the KG-predicted directional effect of the action it
just took against the observed effect.  Accumulated divergence over enough
samples flags a component as dead or inverted; the component's ON and OFF
actions are blacklisted (never \texttt{DO\_NOTHING}, always leaving at
least~1 action), and a \emph{targeted} warm restart zeroes only the
blacklisted actions' Q-values and visit counts before relearning over the
reduced action set.

The \textbf{false-positive guard} is a key design choice: blinds are
interaction-variable-gated (\texttt{Mediates}) actuators whose effect depends
on sunlight, so they are excluded from adjudication; only
unconditional-sign \texttt{Causes} actuators (lamps, spotlight) are judged.
All injected faults are on \texttt{Causes} lamps, so the guard removes a
known false positive without ever masking a designed fault.

Define the headline metric here:
\begin{equation}
  \texttt{RecoveryEpisodes} = \texttt{ReconvergeEpisode}
                             - \texttt{DetectEpisode}
\end{equation}
and the \emph{well-posed-cell} restriction that limits recovery-speed
inference to cells where a post-fault survivor policy actually exists.

\section{Phase~3 --- response-delay learning and KG writeback \evDirect}
\label{sec:approach-phase3}

Phase~3 is an evaluated contribution that does \emph{not} train a Q-table
(\texttt{task\_dynamics.jcm}).  The agent:

\begin{enumerate}[label=(\alph*)]
  \item probes each URI-bearing actuator from a known baseline and measures,
        in simulator ticks, the delay until the controlled zone rank first
        rises;
  \item maintains a numerically stable Welford running mean/variance per
        action and classifies each actuator as instantaneous or delayed
        against a threshold;
  \item writes the learned values back to the KG as a Turtle file with
        \wsRD, \texttt{ws:responseDelayTicks}, sample counts, standard
        deviation, and response class; and
  \item in an exploit phase, scores every action that can reach the target
        zone rank by the tuple $(\textit{infeasible-flag},\
        \textit{energy\_cost},\ \textit{believed\_delay})$ and picks the
        lexicographic minimum.
\end{enumerate}

The KG-primed arm reads the learned delay; the tabula-rasa arm believes
every actuator is instantaneous.  The boundary is explicit \evUnresolved:
richer temporal goal semantics, additional dynamics classes, real-building
deployment, and broader scheduling integration are future work, not claims of
this chapter.

\section{Approach figures}
\label{sec:approach-figures}

Plan three schematic figures.
\begin{description}
  \item[(A) System architecture.]
    BDI agents $\to$ CArtAgO artifacts $\to$ KG/RDF resources $\to$ Node-RED
    simulator $\to$ CSV/log analysis $\to$ evidence ledgers.

  \item[(B) Phase~2 mechanism flow.]
    Clean Q-table warm-load $\to$ unexpected transition $\to$ KG
    expected-effect check $\to$ component blacklist $\to$ warm restart $\to$
    recovery certification.

  \item[(C) Phase~3 pipeline.]
    Slow lab $\to$ actuator probes $\to$ Welford delay table $\to$ \wsRD\
    TTL $\to$ deadline planner $\to$ time-bounded result CSV\@.
\end{description}

\section{Implementation summary (detail in appendix) \evDirect}
\label{sec:approach-impl}

The stack comprises JaCaMo, tabular Q-learner, stereotype reasoner,
\texttt{DynamicsLearner} artifact, Node-RED labs, Python/Node analysis
scripts, and GitHub Actions workflows.  Configuration matrices, exact
commands, and run/artifact provenance are in the appendix.

%%=============================================================
\chapter{Evaluation}
\label{ch:evaluation}
% Target file: content/04_chapter4.tex
%%=============================================================

Structure: experiment design matrix, then per-phase results, then the
statistical-integrity discussion.  Every reported value is tied to a
committed CSV row or a complete derived chain; CIs and bootstrap $p$-values
for Phase~1/2 are flagged as direct analysis outputs that were not
independently resampled.

\section{Experiment design matrix}
\label{sec:eval-design}

Five experiments are evaluated (Table~T-EVAL-1).

\begin{description}
  \item[\texttt{EXP-P1-BUMP} \evDirect, \rsPHE]
    Does bumped clean-lab KG/physics priming beat no-knowledge learning?
    Baseline \qlF, treatment \qlT, profiles \texttt{lab1}/\texttt{lab2}/%
    \texttt{lab3}, selected run \texttt{27461188614}.  Deterministic
    means/differences were recomputed; CI/p/q/Cliff were read from analysis
    CSVs, not re-resampled.  Lab2 is the primary anchor; lab3 reward improves
    more clearly than lab3 binary goal.

  \item[\texttt{EXP-P1-REPL} \evDirect, \rsPHE]
    Do fresh seeds repeat the Phase~1 directions?  Fresh-seed root
    \texttt{P1-BUMP-S11-20}, run \texttt{27462446044}.  This is replication
    evidence, not a new preregistered confirmatory test.

  \item[\texttt{EXP-P1-ABL} \evDirect, \rsPHE]
    Does the targeted cross-zone structure beat equal-budget untargeted
    optimism (\texttt{phase1\_kg\_xzone\_rand})?  Mechanism evidence only;
    language stays exploratory.

  \item[\texttt{EXP-P2-V5} \evDirect, \rsEV; \rsPHE]
    Can clean-trained agents detect faults, discard the defective component,
    and relearn faster with KG knowledge?  Nine fault profiles, modes
    \qlF/\qlT, 10~seeds per CI cell, selected run \texttt{27547019772},
    result commit \texttt{22a448be\ldots}.  Recovery-speed inference
    restricted to \texttt{lab3\_f1dead} and \texttt{lab3\_f1inv} (the only
    well-posed cells).

  \item[\texttt{EXP-P3-DYN-N10} \evDirect, \rsEV; \rsPHE]
    Can the agent learn delays, write them to the KG, and use them for
    deadlines?  Profiles \texttt{lab2\_slow}/\texttt{lab3\_slow}, modes
    \qlT/\qlF, 10~replicas, 6~goals per cell,
    \texttt{seconds\_per\_tick=5.0}, \texttt{blind\_delay\_ticks=12}, run
    \texttt{27621106006}, result commit \texttt{0372ecd5\ldots}.  The n5 run
    is superseded \evSuperseded.
\end{description}

\section{Phase~1 results}
\label{sec:eval-phase1}

\subsection{Lab2 primary anchor \evDerived}

This is the strongest clean-lab result.  \qlT\ $-$ \qlF\ differences:
goal-reaching AUC $+0.0197$ (CI $[0.0147,\,0.0245]$, $q = 0.0$, Cliff's
$\delta = 1.0$); average steps $-0.9$ (CI $[-1.55,\,-0.325]$, $q = 0.0$,
Cliff's $\delta = -0.71$); average redundant actions $-0.9825$ (CI
$[-1.59,\,-0.399]$, $q = 0.0$, Cliff's $\delta = -0.72$); goal rate
$+0.04375$ (CI $[0.0125,\,0.08125]$, $q = 0.009$, Cliff's $\delta = 0.5$).
The directional means and mean differences are deterministic from raw data;
the CIs and $q$-values are original analysis outputs.  KG priming makes lab2
learning faster, more action-efficient, less wasteful, \emph{and} modestly
more successful.

\subsection{Lab2 fresh-seed replication \evDerived}

On fresh seeds, goal AUC difference $+0.0182$ (CI $[0.0106,\,0.0273]$,
$q = 0.0$, Cliff's $\delta = 1.0$) and reward AUC difference $+7.67$
(CI $[2.16,\,12.91]$, $q = 0.033$, Cliff's $\delta = 0.6$).  The direction
holds, supporting robustness of the lab2 finding.

\subsection{Lab3 active boundary \evDerived}

The bumped lab3 \emph{binary-goal} result is \textbf{not} a clean positive:
goal AUC difference $-0.00375$ (CI $[-0.00802,\,0.00067]$, $q = 0.296$),
i.e.\ not distinguishable from zero.  However, the lab3 \emph{reward}
evidence is clearly positive: reward AUC difference $+14.26$ (CI
$[11.23,\,17.20]$, $q = 0.0$, Cliff's $\delta = 1.0$), replicated on fresh
seeds at $+18.85$ (CI $[16.21,\,21.19]$, $q = 0.0$, Cliff's $\delta = 1.0$).
Report binary-goal and reward metrics separately for lab3 and do not
overstate.

\subsection{Superseded as-is lab3 \evSuperseded}

The earlier \enquote{as-is} lab3 rows (goal AUC $-0.00885$, $q = 0.0072$;
cycling $+0.3375$, $q = 0.0252$) are audit/control evidence only, retained
to explain why the bumped targeted run is the lab3 authority.

\subsection{Phase~1 statistical caveat \evUnresolved}

The Phase~1 table has 480~data rows.  Stage~9 confirmed all deterministic
means (348~rows) and mean differences (132~rows) match the raw data with zero
mismatches, but the 348~CI rows and 132~p/q/CI rows were not independently
resampled.  Present Phase~1 CIs/$p$-values as reproduced-from-analysis, not
re-derived.

\section{Phase~2 results}
\label{sec:eval-phase2}

\subsection{Detection and attribution \evDerived}

In the final v5 run, all 18~CI rows have $\texttt{n\_runs} = 10$,
$\texttt{detection\_rate} = 1.0$, and $\texttt{n\_detected} = 10$; and all
18~cells have \texttt{aggregate\_defect\_set\_matches\_expected = true} with
$\texttt{unexpected\_component\_row\_count} = 0$.  In the tested
injected-fault matrix, detection recall is complete and there are no spurious
attributions.  Bound the claim to the tested matrix; do not generalise to
untested fault types or real buildings.

\subsection{Well-posed recovery scope \evDirect}

Only \texttt{lab3\_f1dead} and \texttt{lab3\_f1inv} are marked
$\texttt{well\_posed\_recovery} = \texttt{True}$; all other cells are
descriptive/provisional for recovery-speed inference because a clean
post-fault survivor policy is not well-defined there.

\subsection{Recovery speed \evDerived}

\begin{description}
  \item[\texttt{lab3\_f1dead}: KG is clearly faster.]
    $n_{\text{paired}} = 10$; \qlT\ mean $150.5$ vs.\ \qlF\ mean $363.4$;
    difference $-212.9$ (CI $[-271.9,\,-148.3]$, Wilcoxon $p = 0.0039$,
    Cliff's $\delta = -0.86$, BH $q = 0.0$, \texttt{ql\_true\_faster = True}).

  \item[\texttt{lab3\_f1inv}: KG is \textbf{not} faster.]
    $n_{\text{paired}} = 9$; \qlT\ mean $364.6$ vs.\ \qlF\ mean $332.8$;
    difference $+31.8$ (CI $[-137.2,\,174.6]$, Wilcoxon $p = 0.496$,
    Cliff's $\delta = 0.23$, BH $q = 0.650$, \texttt{ql\_true\_faster = False}).
\end{description}

The supported Phase~2 headline is narrow: KG-primed recovery is faster than
naive for the dead-lamp case but \emph{not} for the inverted-lamp case.

\subsection{Detection latency \evDerived}

Detection timing is mixed across the nine \texttt{DetectEpisode} paired rows
(differences ranging $-2.2\ldots+126.2$); only \texttt{lab3\_f1inv} shows a
large, significant positive detection delay (Wilcoxon $p = 0.00195$, Cliff's
$\delta = 1.0$, BH $q = 0.0$).  Do \textbf{not} state that KG generally
detects faster.

\subsection{Phase~2 statistical caveat \evUnresolved}

Of 521~R2 table rows: 306 are deterministic-from-raw; 99 are
deterministic-or-exact-test; 116 are direct analysis-output resampling rows;
0~failed checks.  Bootstrap CIs and bootstrap $p$-values were not
independently resampled (no Python was available in the audit environment).
The exact Wilcoxon $p$-values, Cliff's $\delta$, and BH~$q$ arithmetic
\emph{were} checked.

\section{Phase~3 results}
\label{sec:eval-phase3}

\subsection{Run authority \evDirect}

Results come from GitHub Actions run \texttt{27621106006} (workflow
\enquote{Phase~3 (process dynamics, response-delay learning)}, success, head
SHA \texttt{3bb5289c\ldots}), published to
\texttt{origin/results@0372ecd5\ldots} under
\texttt{phase3/27621106006-20260616-133306/}.

\subsection{Delay-learning accuracy \evDirect}

Learned slowest (blind) delays stay near the ground-truth 12~ticks:
lab2~\qlF\ $12.1125$, lab2~\qlT\ $12.1875$,
lab3~\qlF\ $12.2125$, lab3~\qlT\ $12.2125$;
maximum relative error $1.77\,\%$.
Across all 180~raw delay rows, the response class matched expectation in
$180/180$ (80~delayed-blind rows, 100~instantaneous non-blind rows).

\subsection{Deadline compliance \evDirect}

\qlF\ compliance: overall $0.5$, tight $0.0$, loose $1.0$ for both slow
profiles.  \qlT\ compliance: $1.0$ on overall, tight, and loose for both
profiles.  The informed planner meets every tight deadline the delay-blind
planner misses.

\subsection{Paired statistics \evDirect}

Overall compliance ($n = 10$ for both profiles): \qlT\ $1.0$ vs.\ \qlF\
$0.5$; difference $0.5$ (CI $[0.5,\,0.5]$, Wilcoxon $p = 0.00195$, Cliff's
$\delta = 1.0$, BH $q = 0.0$).  Tight compliance: difference $1.0$
(CI $[1.0,\,1.0]$, Wilcoxon $p = 0.00195$, Cliff's $\delta = 1.0$, BH
$q = 0.0$).  Loose compliance: difference $0.0$ (both arms already meet loose
deadlines).

\subsection{Energy trade-off \evDirect}

\qlT\ spends more energy: lab2 total-energy difference $+3.2$
(CI $[3.0,\,3.5]$), lab3 $+3.3$ (CI $[3.0,\,3.6]$), with
\texttt{lower\_is\_better = True} and \texttt{ql\_true\_better = False}.
The informed planner reaches for the fast, powered lamp under tight deadlines
instead of the cheap, slow blind.  Lower energy is therefore \emph{not} the
headline success metric---meeting deadlines is---and the energy cost is the
conditioned price of that compliance.

\subsection{Mechanism example \evDirect}

In lab2 rep1, \qlT\ met $6/6$ goals while \qlF\ met $3/6$, missing tight
goals \texttt{g1}, \texttt{g2}, \texttt{g5} after choosing blinds with
believed delay $0.00$.  This single-replica trace makes the aggregate numbers
concrete.

\subsection{Choice-stability nuance \evDerived}

Do not claim actuator-label decisions are bit-identical across replicas:
lab3~\qlT\ loose-goal blind labels split $5/5$ between
\texttt{SetZ1Blinds=ON} and \texttt{SetZ2Blinds=ON} for goals
\texttt{g3}, \texttt{g4}, \texttt{g6}, while every \emph{met} outcome stayed
stable.  The compliance result is robust even though the specific blind chosen
for loose goals is not unique.

\subsection{Superseded n5 run \evSuperseded}

The earlier n5 run \texttt{27598417789} is audit history only: at
$n_{\text{paired}} = 5$ the Wilcoxon $p$ floors at $0.0625$, which cannot
reach significance, so the n10 run is the authority candidate.

\section{Statistical integrity and threats to validity \evUnresolved}
\label{sec:eval-integrity}

Devote an explicit subsection to the evidence limits.  Present the Stage~9
matrix (Figure~F-EVAL-4) separating deterministic/exact checks from
unresolved bootstrap-output categories.  State the threats explicitly:
(i)~no formal preregistration, so all results are exploratory/%
engineering-validation;
(ii)~Phase~1/2 bootstrap CIs and $p$-values are original analysis outputs,
not independently re-resampled;
(iii)~exact workflow-dispatch payloads are \texttt{NOT FOUND} and
artifact-ZIP byte identity is not verified, so committed
\texttt{origin/results} CSV rows are treated as the numeric authority;
(iv)~all evidence is simulator/lab, not real-building deployment.
These are the same items carried into Chapter~5.

\section{Evaluation tables and figures}
\label{sec:eval-figures}

\begin{description}
  \item[T-EVAL-1] Experiment matrix (five experiments, research status,
        baseline, treatment, controls, metrics, key caveat).
  \item[T-EVAL-2] Phase~1 compact table: lab2 anchor, lab2 replication, lab3
        boundary, superseded as-is rows.
  \item[T-EVAL-3] Phase~2 detection/attribution matrix plus well-posed
        recovery table filtered to the two \texttt{lab3~f1} cells.
  \item[T-EVAL-4] Phase~3 delay accuracy, compliance CI, paired compliance,
        energy.
  \item[F-EVAL-1] Phase~1 lab2 forest plot $+$ lab3 boundary panel, marked
        exploratory with the resampling caveat.
  \item[F-EVAL-2] Phase~2 recovery forest plot for the two well-posed cells
        $+$ detection/attribution matrix.
  \item[F-EVAL-3] Phase~3 grouped compliance bars $+$ paired effect panels
        $+$ energy panel.
  \item[F-EVAL-4] Statistical-integrity audit matrix.
\end{description}

%%=============================================================
\chapter{Conclusions}
\label{ch:conclusions}
% Target file: content/05_chapter5.tex
%%=============================================================

\section{RQ answers}
\label{sec:conc-rqs}

\subsection{RQ1 \evInterp}

\textbf{Supported but exploratory.}  The strongest clean-lab evidence is
lab2, where KG/physics priming improves goal AUC, action efficiency,
redundant-action count, and goal rate, with fresh-seed replication.  Lab3 is
split: the reward evidence is positive and replicated, but the binary-goal
metric is not distinguishable from zero.  KG priors accelerate and
streamline clean-lab learning where the lab is well-aligned with the KG,
with the honest caveat that the binary success benefit is lab-dependent.

\subsection{RQ2 \evInterp}

\textbf{Supported as engineering validation.}  The detector and attribution
work across the entire tested injected-fault matrix (complete recall, no
spurious attributions).  The KG \emph{recovery-speed} advantage is narrow
and fault-type-dependent: clearly faster for the dead-lamp case
(\texttt{lab3\_f1dead}), not faster for the inverted-lamp case
(\texttt{lab3\_f1inv}).  The capability claim (detect $\to$ discard $\to$
relearn works) must be stated separately from the weaker speed claim.

\subsection{RQ3 \evInterp}

\textbf{Supported as engineering validation.}  The agent learns the blind's
response delay accurately (${\leq}1.77\,\%$ relative error), writes it back
to the KG, and the informed planner meets every tight deadline the
delay-blind planner misses---at a consistent, explainable energy cost.  Bound
this to the two tested slow-lab profiles with a single delayed actuator class.

\section{Limitations \evUnresolved}
\label{sec:conc-limitations}

The following limitations are stated without softening.

\begin{enumerate}
  \item No formal post-pivot preregistration was found; no result is
        preregistered confirmatory.
  \item Phase~1/2 bootstrap CIs and bootstrap $p$-values were not
        independently resampled.
  \item Exact workflow-dispatch input payloads are \texttt{NOT FOUND} for all
        phases.
  \item Local artifact directories are not ZIP-byte-verified against GitHub
        artifact digests.
  \item All evidence is from controlled simulators, not deployed buildings.
  \item The lab3 stale-magnitude documentation issue remains unresolved.
  \item DOI-only related-work rows are not full-text reviewed (GAP-046).
  \item Phase~3 covers a limited dynamics scope: one delayed actuator class
        in two profiles.
\end{enumerate}

\section{Future work \evInterp}
\label{sec:conc-future}

Natural continuations include: more actuator/dynamics classes and richer
temporal goal specifications; online/real-building deployment;
artifact-ZIP verification and recovery of exact dispatch payloads;
preregistered confirmatory reruns; independent bootstrap reruns to close the
resampling gap; and a full-text related-work audit of the DOI-only sources.

%%=============================================================
\appendix
\chapter{Source and Run Provenance}
\label{app:provenance}
% Target file: content/appendix.tex  Appendix A
%%=============================================================

Record the repository state during Stage~11: branch
\texttt{phase3-process-dynamics}, HEAD \texttt{eda6ca1f\ldots},
\texttt{origin/results} \texttt{0372ecd5\ldots}.

\begin{description}
  \item[Phase~1 route.]
    Stage~6 \texttt{PHASE1\_TABLES.csv} is the authoritative numeric table
    (480~rows, zero deterministic mismatches).

  \item[Phase~2 route.]
    Run \texttt{27547019772}, v5 head \texttt{985c7a18\ldots}, result commit
    \texttt{22a448be\ldots}.

  \item[Phase~3 route.]
    Run \texttt{27621106006}, head \texttt{3bb5289c\ldots}, artifact
    \texttt{phase3-consolidated} id \texttt{7668397627}, digest
    \texttt{sha256:a882677d\ldots}, result commit \texttt{0372ecd5\ldots}.
\end{description}

%%=============================================================
\chapter{Reproduction Commands}
\label{app:reproduction}
% Appendix B
%%=============================================================

Commands below are recorded as notes, not claims that they were rerun in
Stage~11.

\begin{itemize}
  \item Phase~1 table extractor:
        \texttt{node paper\_notes/phase1\_extract\_tables.mjs \ldots}
  \item Phase~2 extractor: (see \texttt{APPENDIX\_NOTES.md} for exact
        invocation).
  \item Stage~9 audit script: (see \texttt{APPENDIX\_NOTES.md}).
  \item Phase~3 n10 recompute:
        \texttt{python analysis/phase3\_dynamics.py --root \ldots{} --out
        \ldots}
  \item Tree-hash command: (see \texttt{APPENDIX\_NOTES.md}).
  \item Phase~3 n5 recompute: \evSuperseded\ (n5 run superseded by n10).
\end{itemize}

%%=============================================================
\chapter{Experiment Matrices}
\label{app:matrices}
% Appendix C
%%=============================================================

\begin{description}
  \item[Phase~1.]
    Profiles \texttt{lab1}/\texttt{lab2}/\texttt{lab3}; stereotype arms
    true/false; benchmark modes \texttt{rule\_based}/\qlF/\qlT; component-%
    isolation run modes \texttt{phase1\_baseline}\ldots%
    \texttt{phase1\_kg\_xzone\_rand}.

  \item[Phase~2.]
    Nine fault profiles $\times$ two modes; parent ports, flows, suffixes,
    and state dimensions documented in \texttt{APPENDIX\_NOTES.md}.

  \item[Phase~3.]
    Profiles \texttt{lab2\_slow}/\texttt{lab3\_slow}; two modes;
    \texttt{seconds\_per\_tick = 5.0}; \texttt{blind\_delay\_ticks = 12};
    \texttt{probes\_per\_actuator = 8}; settle/poll/max-wait constants;
    target rank~3; probe sun~900; six goals.
\end{description}

Exact dispatch inputs remain \texttt{NOT FOUND} for all phases.

%%=============================================================
\chapter{Statistical Integrity}
\label{app:stats}
% Appendix D
%%=============================================================

Reproduce the Stage~9 audit breakdown.

\begin{description}
  \item[Phase~1 (480 rows).]
    348~mean rows + 132~diff rows verified (deterministic); 348~CI rows +
    132~p/q/CI rows not independently resampled.

  \item[Phase~2 (521 rows).]
    306~deterministic-from-raw; 99~deterministic-or-exact-test;
    116~direct-analysis-output; 0~failures.
    \texttt{RecoveryEpisodes} family size~2;
    \texttt{DetectEpisode} family size~9.

  \item[Phase~3 recompute.]
    4~delay rows; 4~compliance CI rows; 8~paired rows; delay and paired CSVs
    byte-identical to the downloaded consolidated CSVs; compliance CI a
    numeric (not byte) match.
\end{description}

BH nuance: BH arithmetic was checked from aggregate bootstrap $p$-values,
but those bootstrap $p$-values are themselves direct analysis outputs.

%%=============================================================
\chapter{Full Table Routes}
\label{app:tables}
% Appendix E
%%=============================================================

\begin{itemize}
  \item Phase~1 compact evidence table: \texttt{R2P1T-001}\ldots%
        \texttt{R2P1T-014}.
  \item Phase~2 final CI table: 18~rows.
  \item Phase~2 final paired table: rows~2--3 \texttt{RecoveryEpisodes};
        rows~4--12 \texttt{DetectEpisode}.
  \item Phase~2 detailed verification table: 522~rows with input chains.
  \item Phase~3 result tables: see \texttt{PHASE3\_TABLES.csv}.
  \item Bibliographic ledger: \texttt{BIBLIOGRAPHIC\_LEDGER.csv},
        \texttt{BIB-S10B-001}\ldots\texttt{016} + 1~rejected entry.
\end{itemize}

%%=============================================================
\chapter{Figure Specifications}
\label{app:figures}
% Appendix F
%%=============================================================

Route the architecture component/dataflow diagrams, the Phase~1/2/3 figure
families, and the related-work map as evidence-grade figures (not yet
rendered).  Full specifications are in \texttt{APPENDIX\_NOTES.md}.

%%=============================================================
\chapter{Audit History Index}
\label{app:audit}
% Appendix G
%%=============================================================

\begin{description}
  \item[Phase~1.] as-is $\to$ bumped $\to$ ablation supersession chain.
  \item[Phase~2.] v1--v5 (v5 removes the earlier \texttt{SetSpotlight}
        attribution rows).
  \item[Phase~3.] n5 $\to$ n10 (n5 had the Wilcoxon-floor issue).
  \item[Bibliographic.] Rejected Brick DOI correction
        (\texttt{BIB-S10B-REJ-001} $\to$ \texttt{BIB-S10B-010}).
\end{description}

%%=============================================================
\chapter{Evidence-Gap Index}
\label{app:gaps}
% Appendix H
%%=============================================================

\begin{description}
  \item[GAP-001] Advisor source provenance.
  \item[GAP-021] Formal preregistration: \texttt{NOT FOUND}.
  \item[GAP-035] Phase~1 result limits.
  \item[GAP-037] Phase~2 result limits.
  \item[GAP-044] Stage~9 statistical boundary.
  \item[GAP-045] Phase~3 authority limits.
  \item[GAP-046] Related-work full-text review (DOI-only rows).
  \item[GAP-047] (see active ledger).
\end{description}

%%=============================================================
%% Cross-cutting discipline reminders (comment block, not typeset)
%%
%% - Never upgrade exploratory/engineering-validation to confirmatory.
%% - Separate deterministic numbers from not-independently-resampled CIs.
%% - Keep [SUPERSEDED] material in audit history only.
%% - Use NOT FOUND / NOT VERIFIED explicitly; no gap-filling by inference.
%% - Bound every empirical claim to the simulator/lab + tested matrix.
%%=============================================================

\end{document}
